\documentclass[11pt, a4paper]{article}

\usepackage[
    top=2.3cm, bottom=2.3cm,
    left=2.5cm, right=2.5cm
]{geometry}
\usepackage{microtype}
\usepackage{xcolor}
\usepackage{enumitem}
\usepackage{hyperref}
\hypersetup{colorlinks=true, urlcolor=blue, linkcolor=black}

% A reviewer comment, followed by our reply.
\newcommand{\point}[1]{\medskip\noindent\textit{#1}\par}
\newcommand{\reply}{\textbf{Response.}\ }

\setlength{\parindent}{0pt}
\setlength{\parskip}{7pt}

\begin{document}

\begin{flushright}
    Professor Xinzhong Zhu \quad (corresponding author) \\
    Moatassim Barakat \\
    College of Computer Science and Technology \\
    Zhejiang Normal University, Jinhua 321004, China \\
    \href{mailto:zxz@zjnu.edu.cn}{zxz@zjnu.edu.cn}
    \quad
    \href{mailto:moatassimbarakat@gmail.com}{moatassimbarakat@gmail.com} \\[4pt]
    \today
\end{flushright}

\bigskip

Professor Tamas Sziranyi \\
Area Editor, \textit{Digital Signal Processing}

\bigskip

Manuscript DSP-D-26-02086 \\
Original title: ``MEISCF: Multi-scale Edge Information Selection with
Cross-scale Fusion for Small Object Detection in UAV Imagery''

\bigskip

\begin{center}
\large\textbf{Response to Reviewers}
\end{center}

\bigskip

Dear Professor Sziranyi,

Thank you for the opportunity to revise this manuscript, and please pass our
thanks to the six reviewers. Their central request was the same one, put in
different ways: isolate the contribution of each component under matched
conditions, and show that the reported gains are larger than the noise.

We ran those controls. They did not confirm the paper. Under a matched epoch
budget, almost all of the accuracy we had attributed to our modules came from the
training schedule, and the direction of that effect was the opposite of what we
had claimed: single-resolution training at the deployment size beat our
four-phase progressive schedule by 3.3 to 4.2 points mAP@50. A second seed
reversed the sign of the MEIS and FRM difference. A factorial ablation showed
that neither multi-scale dilation nor adaptive gating, the two ingredients that
define MEIS, did better than a plain convolution with the same capacity. Resolved
by object size, the modules improved no size range below 32 pixels.

We could not honestly repair the original manuscript, so we rewrote it as the
controlled study the reviewers asked for. The revised paper reports what the
experiments show, including the results that refute our earlier claims. Its new
title is:

\begin{quote}
``What Drives Small Object Detection Accuracy in UAV Imagery? A Controlled Study
of Training Resolution, Detection Stride, and Feature Modules''
\end{quote}

We recognise that this is a larger change than a revision normally involves, and
that the paper now argues against the method it was named after. We think it is
the more useful paper, and the one the reviewers' comments point to. If the
Editor judges that it is better handled as a new submission, we will follow that
advice.

\section*{Summary of the main changes}

\begin{itemize}[leftmargin=1.3em, itemsep=2pt]
\item \textbf{Matched comparisons.} Every comparison now varies one factor at a
time at a matched epoch budget, with training and evaluation at the same input
size. The schedule is crossed with the detection head in a $2\times2$ design
(Table~5).
\item \textbf{Statistical testing.} All differences are tested with a paired
bootstrap and a paired permutation test over the 548 validation images, with
Bonferroni correction across the ten classes. Two seeds are reported for the
module comparison and for both schedules of the four-level detector
(Sections~4.5, 5.4).
\item \textbf{Component isolation.} A factorial ablation separates multi-scale
dilation from adaptive gating and compares both against a plain convolution of
equal capacity (Section~5.3.1, Table~7).
\item \textbf{Accuracy by object size.} Results are resolved into COCO size
ranges and into bands of network-input pixels tied to the detection strides
(Section~5.5, Table~8, Fig.~8).
\item \textbf{A predictive account.} Section~3 states a sampling argument and
derives three predictions before the experiments, so the architectural results
can be judged against a stated expectation rather than explained afterwards.
\item \textbf{Cross-dataset evaluation.} UAVDT is now evaluated zero-shot with
ignore regions honoured and sequence-level resampling, and the classes it cannot
test are named (Section~5.6, Table~9).
\item \textbf{Cost.} Throughput is measured where the GPU, not the host, is the
limit, and reported next to the operation counts (Section~4.2, Table~2).
\item \textbf{Claims removed.} The comparison against CF-YOLO under unmatched
training, the ``+17.6\%'' and ``+10.2\%'' headline numbers, the claim of novelty
for three-level fusion, and the unsupported activation-magnitude statistic have
all been removed.
\end{itemize}

\textit{A note on cross-references.} The manuscript was rewritten rather than
edited, so line numbers of the original no longer locate anything. We therefore
point to sections, tables and figures of the revised manuscript.

% ===========================================================================
\section*{Reviewer \#1}

\point{Introduction section could more clearly distinguish the genuine
methodological contribution from the integration of previously established
techniques. Several components \dots are based on well-known concepts. Please
explicitly identify which aspects constitute the primary scientific
contribution.}

\reply We agree, and the controls we ran to answer this question showed that the
components did not carry the result. The contribution is no longer a module. It
is the controlled comparison itself: which of three design decisions --- training
resolution schedule, detection stride, and feature modules --- actually moves
small-object accuracy, measured one factor at a time with paired tests and two
seeds, and resolved by object size. The Introduction now states this and lists
three contributions: the matched-budget schedule-by-stride comparison, the
controlled evaluation of the modules including the factorial ablation, and the
sampling account with its size-resolved test (Section~1). We also state plainly
that the modules we had proposed do not survive the controls.

\point{The Related Work section \dots is largely descriptive. I recommend more
critical discussion of previous approaches and a clearer explanation of how the
proposed framework differs from recent lightweight YOLO variants,
Transformer-based detectors, and other state-of-the-art methods.}

\reply Related Work has been restructured around the four questions the study
tests: feature modules for small aerial objects, resolution and sampling,
resolution schedules during training, and variance and significance in detector
evaluation (Sections~2.1--2.4). Each subsection now asks what evidence the cited
work provides rather than what it proposes: how many seeds, whether the
comparison was matched, and whether the reported gain exceeds the run-to-run
spread. Section~2.4 is new and reviews what is known about seed variance in
detection benchmarks, which is the standard against which we judge our own
effects.

\point{The largest improvement is achieved by the progressive multi-resolution
training strategy, whereas the architectural modifications contribute only an
additional improvement \dots a more convincing evaluation should include
comparisons with competing approaches, particularly CF-YOLO, trained using the
same progressive training schedule.}

\reply Your reading was correct, and the controlled experiment went further than
we expected: the schedule dominates, and it dominates in the opposite direction.
Single-phase training at 1280\,px beat the four-phase progressive schedule by
3.94 points mAP@50 with seed 0 and 4.16 with seed 42, with and without the
stride-4 level (Table~4, Table~5; paired bootstrap 95\% CI $[4.1,\,5.5]$ at
1536\,px, $p<0.002$, permutation $p=0.001$). At matched training compute the
single-phase run was 3.87 points ahead (Section~5.1, Fig.~4).

We did not retrain CF-YOLO under our schedule. The comparison is no longer
load-bearing, because the revised paper makes no claim of superiority over
CF-YOLO or any other published method; that claim and its headline number have
been removed. Section~5.7 and Table~10 now present published results with their
protocols side by side --- split, input size, initialization --- precisely to show
that they are not a ranking. Retraining a third-party detector under two
schedules would be a worthwhile experiment, but its outcome would not change any
conclusion we now draw, and we would rather not publish claims about another
group's method from a reimplementation we cannot validate.

\point{Adaptive gating mechanism within the MEIS module and the selected dilation
rates are evaluated only as part of the complete framework.}

\reply This is now a dedicated factorial experiment (Section~5.3.1, Table~7,
Fig.~7). We crossed the dilation rates, $(1,2,4)$ against $(1,1,1)$, with the gate
type, input-adaptive against a static learnable channel weight, holding the
parameter count fixed. The corner with $(1,1,1)$ dilations and a static gate is a
plain residual block of depthwise-separable convolutions with the same capacity as
MEIS. Neither ingredient beat that plain block on two seeds. Whatever MEIS
contributes, it is not multi-scale dilation and not adaptive gating.

\point{The methodological novelty of FRM also appears relatively limited \dots
its structure is essentially equivalent to the well-known CBAM module, while the
primary difference lies only in its placement \dots I could be taking into account
deeper theoretical analysis explaining why this particular combination of modules
should produce more discriminative feature representations.}

\reply We accept the criticism and have stopped claiming otherwise. FRM is
described as CBAM applied after fusion, with a citation to CBAM and no claim of
structural novelty (Section~4.2). It is no longer counted as a separate
contribution with its own percentage.

On the request for a theoretical account: rather than argue after the fact why
the combination should work, Section~3 now states a sampling argument
\emph{before} the experiments and derives three testable predictions from it, with
the object-size distribution of VisDrone measured in network-input pixels
(Table~1, Figs.~1--2). The account predicts the stride results correctly and fails
on the schedule result, and we say so (Section~6.1). This seemed to us more useful
than a post-hoc justification of a combination that the measurements do not
support.

\point{Another aspect that deserves further attention is the interpretability
\dots I recommend present visualizations of feature maps, activation maps, or
attention maps demonstrating that the network indeed focuses on edge-related
structures.}

\reply We have not added these, and we would like to explain why rather than
simply decline. Such a visualization would be evidence for a claim the paper no
longer makes. It can show that a branch responds to boundaries, but not that this
response contributes to detection accuracy, and our measurements say it does not:
MEIS improves no size range below 32 pixels on either seed, and a plain
convolution of equal capacity does as well (Section~5.3, Table~8). Publishing
attention maps beside that result would invite the reader to trust the picture
over the measurement. We report the quantitative decomposition instead, by object
size and by component.

\point{The cross-dataset evaluation on the UAVDT dataset considers only the
dominant car category and results in a relatively modest improvement. Please \dots
evaluate the robustness \dots under challenging imaging conditions, such as motion
blur, varying illumination, partial occlusion, and image noise.}

\reply The UAVDT evaluation has been rebuilt: zero-shot at 1280\,px, official
ignore regions honoured, vans merged into cars to match UAVDT's label set, and
confidence intervals from resampling whole sequences rather than frames, since
frames within a sequence are not independent (Section~5.6, Table~9). The car
ordering reproduces VisDrone --- schedule $+2.24$ and $+2.47$ AP, stride $+1.82$,
MEIS $-0.81$ --- but no sequence-level interval excludes zero, and we now report
that rather than the point estimates alone. We also state why trucks and buses
cannot carry a conclusion: each has 62\% of its boxes in a single sequence, so
their AP moves by tens of points between resamples and the two classes rank the
models in opposite orders.

We did not run a controlled corruption study. Doing it properly means calibrated
blur, noise and illumination levels applied to a fixed image set, which is a
separate experiment of its own size; and since the modules show no benefit in
clean conditions on either seed, a robustness study would most likely measure the
robustness of the baseline. We list it as future work in Section~6.4 rather than
report a partial version of it.

\point{I recommend extending the experimental comparison to include more recent
lightweight Transformer-based object detectors as well as the latest YOLO variants
specifically designed for aerial imagery.}

\reply Section~2.1 now discusses a recent detection transformer for this task,
UAV-DETR (Zhang et al., arXiv:2501.01855), alongside the aerial YOLO variants, and
Section~2.4 was added to review how this literature handles run-to-run variance.

We did not, however, add new detector families to the experiments, for a reason
that is central to the paper: comparing architectures without matched schedules,
matched input sizes, matched epoch budgets and paired tests would reproduce
exactly the problem the paper documents. A schedule change alone moves mAP@50 by
about 4 points in our measurements, which is of the same order as the gains those
detectors report over their own baselines --- UAV-DETR reports $+4.2$ points AP@50
on VisDrone --- so a table of detectors trained under their own recipes cannot
separate architecture from protocol. Section~5.7 therefore reports published
numbers with their protocols attached and explicitly declines to rank them, and
Section~6.4 lists the narrow architectural scope --- one detector family, one
backbone scale --- as a limitation of this study.

% ===========================================================================
\section*{Reviewer \#2}

\point{\dots the manuscript is clearly written, the technical contribution is
well supported by the experimental results, and the proposed approach represents a
valuable contribution to the literature.}

\reply We thank the reviewer for the positive assessment, and we owe them a
straightforward explanation of why the manuscript they recommended has changed so
much. The controls requested by the other reviewers did not support the
framework's effectiveness. The experimental design the reviewer found
comprehensive has been kept and extended, but its conclusions have reversed: the
gains we reported came from the training schedule rather than from the modules,
and the schedule itself works in the opposite direction to the one we claimed. We
hope the revised paper remains useful to the community for the measurement
practice it documents, even though its findings are negative for the method.

% ===========================================================================
\section*{Reviewer \#3}

\point{The motivation is not clear. There are several methods combined together,
what is the novelty of this manuscript and what are their functions?}

\reply The manuscript no longer combines methods. It tests three design decisions
separately and reports which one matters. The motivation, stated at the start of
Section~1, is that gains of one to three points are routinely attributed to new
modules in this literature while the training protocol is left uncontrolled and
run-to-run variance is not reported; our own earlier results were an instance of
this. The novelty is the controlled decomposition and its outcome, not a module.

\point{Figs.~2 and 3 are confusing.}

\reply The figure set has been rebuilt. The revised manuscript has eight figures:
the sampling geometry that motivates the study, showing the same object under
three stride and input-size settings (Fig.~1), the class distribution of the
training set (Fig.~2), the detector
architecture with the stride-4 path marked (Fig.~3), accuracy against training
compute (Fig.~4), per-class differences (Fig.~5), a forest plot of every headline
contrast with its confidence interval (Fig.~6), the factorial ablation (Fig.~7),
and accuracy by object size (Fig.~8). Each is drawn to make one comparison, and
the caption states what it shows.

\point{The expression should be more concise. The contributions should be
shortened.}

\reply The manuscript has been rewritten throughout. The contribution list is now
three short items in Section~1 rather than an extended enumeration, and the
abstract is a single paragraph of results.

\point{Please check the parts seem like being written by AI, for example, the
former part of the abstract.}

\reply The abstract has been written from scratch and now opens with the question
the study asks and the design used to answer it, with no promotional framing. The
manuscript as a whole was rewritten in plain declarative prose. Our use of
AI-assisted tools is declared through the submission system, in line with the
journal's policy.

\point{The references should be logically organized with enough analyses, and
some outdated references should be deleted or replaced. Mohamed Meyer Kana Kone,
et al. Applied Soft Computing, 2026, 203: 116069.}

\reply Related Work is now organized by the four questions the study tests rather
than by method family, and each entry is discussed in terms of the evidence it
provides (Sections~2.1--2.4). The reference list was rebuilt for the new
manuscript, superseded entries were removed, and the author lists and page ranges
of the remaining entries were checked against the publishers' records.

We were not able to locate the suggested Applied Soft Computing article from the
details given, and the work we could find under that author name concerns plant
disease identification and image watermarking rather than object detection in
aerial imagery. We would gladly add it if the reviewer can supply the title, and
we would equally welcome pointers to work that bears on the four questions above
--- particularly any aerial-detection study that reports run-to-run variance,
since Section~2.4 had to draw most of its evidence from outside this literature.

\point{Can quantum computing be helpful to solve this problem? NR Zhou, et al.
\dots Gong LH, et al. \dots}

\reply It is an interesting question, but outside what this manuscript can speak
to. The study is an empirical comparison of training schedules, detection strides
and feature modules for a convolutional detector on classical hardware, and none
of its measurements bear on quantum methods. We have not added these citations,
since we could not connect them to any claim in the paper, and a citation we
cannot justify in the text would not serve the reader.

\point{It is hard to follow since the theoretical foundation is not clear.}

\reply Section~3 is new and supplies it. It sets out how the detection stride
determines the sampling density of an object of a given size, measures the
VisDrone size distribution in network-input pixels at each training resolution
(Table~1), and derives three predictions: which size ranges a stride-4 level
should help, which it should hurt, and what a module acting after downsampling can
and cannot recover. The experiments then test those predictions, and Section~6.1
reports which held and which failed.

\point{Why should we use YOLOv11 but not other new versions?}

\reply Section~4.2 gives the reason: the study needs a detector whose neck and
head can be modified in a controlled way and whose training pipeline is public and
stable, so that the only difference between two runs is the factor under test.
YOLOv11s provides this with COCO-pretrained weights at a scale that fits the
repeated multi-seed training the design requires. The conclusions concern
schedule, stride and module placement, which are not specific to the version;
Section~6.4 states the single-family scope as a limitation.

\point{More comparisons would be better to convince the audience. The results
should be explained theoretically.}

\reply The revised paper adds the comparisons that isolate causes rather than
lengthen the ranking: the $2\times2$ schedule-by-stride design (Table~5), the
factorial ablation of MEIS (Table~7), size-resolved accuracy for every contrast
(Table~8), a second seed for four configurations, and a zero-shot cross-dataset
evaluation (Table~9). Each result is set against the prediction of Section~3, and
Section~6.1 separates what the sampling account explains --- the whole pattern of
the stride results --- from what it does not --- the schedule effect, for which we
give the competing explanations and say which experiment would decide between
them.

\point{What is the cost of the performance improvement?}

\reply Cost is now reported for both kinds of change. For training, we report
pixel-epochs, a hardware-independent measure of the pixels a schedule processes
(Section~4.3, Table~4), and accuracy against compute (Fig.~4). For inference,
Table~2 reports parameters, operation counts at three input sizes and measured
throughput at 1280 and 1536\,px. The stride-4 level costs 74\% more operations
and 89\% more time; MEIS with FRM costs 18\% more operations but 85\% more time,
so its real cost is 1.6 times what its operation count implies, for no measurable
accuracy benefit.

\point{The advantages and disadvantages of this designed method were detailed, can
it be realized with current technology?}

\reply All models are ordinary convolutional detectors that train and run on a
single consumer GPU; every number in the paper was produced on one RTX~3080, and
Table~2 gives the measured throughput. The revised manuscript is explicit that the
modules' disadvantages outweigh their advantages: they cut throughput nearly in
half at 1280\,px and improve no size range below 32 pixels.

% ===========================================================================
\section*{Reviewer \#4}

\point{1. The +1.6\% gain from the MEIS module is a combined result of three
design dimensions \dots Is it the ``edge extraction'' itself that is effective? Or
does the ``adaptive gating'' contribute the gain? Or is it simply that after
resolution enhancement, any additional convolution can bring a +1.x\%
improvement?}

\reply The third possibility is the closest to what we measured. The factorial
ablation (Section~5.3.1, Table~7) crosses dilation rates with gate type at a fixed
parameter count, and its plain-convolution corner --- $(1,1,1)$ dilations, static
gate --- performs as well as full MEIS on two seeds. Moreover the $+1.6\%$ itself
does not survive a second seed: measured against Baseline-P2 under matched
conditions, MEIS with FRM changed mAP@50 by $+0.58$ points with seed 0 and
$-0.34$ with seed 42 at 1536\,px, with paired bootstrap intervals of
$[-0.15,\,1.27]$ and $[-1.00,\,0.33]$ and permutation tests that agree
(Section~5.3, Table~6, Fig.~6). Even on the favourable seed the interval excludes
improvements larger than about 1.3 points. Split by object size, COCO small
objects lost 0.8 points averaged over the two seeds, 95\% CI $[-1.5,\,-0.1]$
(Section~5.5).

\point{2. The authors claim to ``simultaneously integrate P3/P4/P5,'' but in fact
only the P4 level achieves three-layer fusion \dots Moreover, learnable scalar
weights have already been widely used in CF-YOLO and EfficientDet.}

\reply Both observations are correct and the claim has been withdrawn. The neck is
now described accurately, with its normalized learnable weighting attributed to
the fast normalized fusion of BiFPN (Section~4.2), and no novelty is claimed for
it. It is not part of the main experiments; the runs that used it are reported in
Appendix~A as earlier single-seed ablations, clearly labelled as such.

\point{3. The authors acknowledge that CF-YOLO did not use progressive training
and that the comparison is an ``upper bound.'' However, this comparison is
repeatedly emphasized as a +10.2\% improvement in the Abstract, Introduction, and
Conclusion, which can easily mislead readers.}

\reply Agreed, and it has been removed from the abstract, the introduction, the
conclusion and everywhere else. The revised manuscript makes no comparative claim
against CF-YOLO. Section~5.7 explains why published VisDrone numbers are not
comparable across papers --- split, input size, initialization and even the AP
definition differ --- and Table~10 lists protocols alongside numbers so that the
rows cannot be read as a ranking.

\point{4. On UAVDT, evaluation is conducted only on the ``car'' class \dots this
also reveals that MEISCF's cross-domain generalization capability for truly small
objects has not been validated at all.}

\reply Correct, and this is now stated in the abstract rather than left to the
discussion: UAVDT contains only vehicle classes, so cross-domain generalization
for small non-vehicle objects remains untested. The rebuilt evaluation
(Section~5.6) reports what the benchmark can and cannot support: cars reproduce
the VisDrone ordering but no sequence-level interval excludes zero; trucks appear
in six of eleven sequences with 62\% of them in one, and buses in four sequences
with 62\% in one, so neither class can carry a conclusion. UAVDT's objects are also
larger, with 30.7\% of instances below 32 pixels against 68.6\% on the VisDrone
validation set (Section~4.1), so the two benchmarks do not probe the same size
regime.

\point{5. Table~1 and Table~3 appear to be redundant \dots However, the reported
metrics for the same methods are inconsistent across the two tables. Why?}

\reply The inconsistency was real and we apologize for it: those tables mixed runs
from different configurations and input sizes. The revised manuscript uses one
evaluation protocol for every number it reports --- letterbox to the stated size,
confidence 0.001, NMS IoU 0.5, at most 600 detections per image (Section~4.4) ---
and any number produced under a different configuration is confined to
Appendix~A, which states the four ways that configuration differed. We also
identified a second source of inconsistency that affects comparisons with the
literature: Ultralytics releases differ in how the precision--recall curve is
terminated, which changes mAP@50 by 0.5 to 0.7 points on identical detections.
Section~4.4 states the definition we use.

\point{6. Some experiments are conducted at 1024$\times$1024 resolution, while
others are at 1280$\times$1280. What is the rationale behind this choice? Why is
the resolution not kept consistent across all experiments?}

\reply There was no good rationale; the resolutions reflected the order in which
the experiments had been run. In the revised manuscript every model is trained at
its stated size and evaluated at 1280 and 1536\,px, the same two sizes for all
models. The 1024\,px runs appear only in Appendix~A, and one further 1024\,px
evaluation is used deliberately in Section~5.5 to test whether the schedule effect
is a train--test scale mismatch, which it is not.

% ===========================================================================
\section*{Reviewer \#5}

\point{1. The authors acknowledge that the individual contribution of the
per-input channel-wise gating \dots was not isolated due to hardware constraints
\dots the authors should consider adding a lightweight theoretical justification
or a highly scaled-down ablation \dots}

\reply We took this suggestion and were able to run the full ablation rather than
a scaled-down one, using a shortened progressive schedule and two seeds
(Section~5.3.1, Table~7, Fig.~7). Gate type and dilation rates are crossed at a
fixed parameter count. Neither ingredient beat a plain convolution of equal
capacity, so the $+1.6\%$ attributed to MEIS cannot be assigned to either of the
mechanisms that define it --- and under matched conditions with two seeds, the
$+1.6\%$ itself does not replicate.

\point{2. The cross-dataset evaluation on UAVDT focuses exclusively on the ``car''
category \dots the authors should explicitly state in the abstract that
``cross-domain generalization for small non-vehicle objects remains unvalidated.''}

\reply Done, in those terms. The abstract now ends: ``UAVDT has only vehicle
classes, so small non-vehicle generalization remains untested.''

\point{3. The +10.2\% improvement over CF-YOLO is rightly noted as an upper bound
\dots A brief theoretical discussion \dots estimating how much of the +12.8\%
curriculum learning gain might transfer to CF-YOLO would strengthen this
comparative analysis.}

\reply The premise of the estimate no longer holds: the curriculum did not
produce a gain. Trained at a matched epoch budget, the progressive schedule was
3.3 to 4.2 points \emph{worse} than single-resolution training at the deployment
size, on both detection heads and both seeds (Tables~4 and~5). The comparison
against CF-YOLO has been removed together with the claim it supported, so there is
no transfer estimate to make.

\point{4. To strengthen the literature review \dots The authors are encouraged to
review and integrate the following relevant frameworks: (a) UAV-DETR \dots (b)
PADrone: Pre-Flight Abnormalities Detection on Drones Using Deep RF Sensing \dots
(c) Multi-Sensor Fusion for UAV Classification Based on Feature Maps of Image and
Radar Data.}

\reply We have restructured Related Work around the four questions the study tests
(Sections~2.1--2.4), and we have added a detection transformer to Section~2.1,
which also answers Reviewer~\#1's request for a recent Transformer-based
detector.

One clarification on suggestion~(a). Two different papers share the name
UAV-DETR. The one at the identifier given, arXiv:2603.22841, is
\textit{UAV-DETR: DETR for Anti-Drone Target Detection}, which detects drones as
targets --- a different task from ours, where the camera is on the drone and the
targets are small ground objects. The one that shares our task and our benchmark
is Zhang et al., \textit{UAV-DETR: Efficient End-to-End Object Detection for
Unmanned Aerial Vehicle Imagery} (arXiv:2501.01855), which adds frequency-domain
feature fusion, detail-preserving downsampling and semantic alignment to RT-DETR
and reports $+4.2$ points AP@50 on VisDrone. We cite that one, and we use it to
make a point the study depends on: its reported gain is of the same order as the
effect we measure from the training schedule alone, which is why we compare
factors rather than architectures.

On suggestions~(b) and~(c): this manuscript studies single-sensor optical
detection of objects in aerial images, and its measurements concern detection
stride, training resolution and module placement. The RF-sensing and image--radar
fusion works address a different sensing modality and a different task, drone
classification and pre-flight anomaly detection rather than detection of objects
within an image, so we could not connect them to a claim in the paper without
padding the review.

% ===========================================================================
\section*{Reviewer \#7}

\point{1. The paper attributes the +17.6\% total gain to a combination of training
and architecture, but the +12.8\% from progressive training is measured against a
baseline that is not directly comparable. The architectural modules are only
evaluated on top of progressive training, so their contribution is not isolated
from the benefits of higher resolution inputs.}

\reply This diagnosis was right, and following it changed the paper's conclusions.
The revised design controls the training schedule explicitly: the progressive
schedule and a single-phase schedule at 1280\,px are crossed with the three-level
and four-level heads at a matched budget of up to 350 epochs (Table~5). With the
comparison made fair, the progressive schedule loses by 3.3 to 4.2 points mAP@50.
The modules are evaluated against Baseline-P2 under an identical schedule and
identical inputs, and the $+17.6\%$ figure has been withdrawn.

\point{2. The paper admits that FRM is ``a placement choice rather than a
structural novelty'' \dots However, the paper still counts FRM as a distinct
module contributing +1.3\%. This is a significant overstatement of novelty.}

\reply Agreed and corrected. FRM is presented as CBAM applied after fusion, with
no claim of novelty and no separate percentage attributed to it (Section~4.2). Its
effect is measured only as part of the MEIS\,+\,FRM contrast, which is not
significant on either seed.

\point{3. \dots the +10.2\% gap ``represents an upper bound under unmatched
training conditions.'' This comparison is misleading and should not be presented
as a fair state-of-the-art comparison.}

\reply Removed. See our response to Reviewer~\#4, point~3.

\point{4. Curriculum learning through progressive resolution increase is a
well-established technique \dots The paper presents this as a key innovation, but
it is a standard training trick, not a methodological novelty.}

\reply We no longer present it as an innovation. Section~2.3 places progressive
resolution training in its literature, and the revised paper treats it as a factor
to be tested rather than a contribution to be claimed --- a test it does not pass,
since it underperforms single-resolution training at the deployment size under a
matched epoch budget.

\point{5. The paper reports a +25\% increase in GFLOPs \dots and a +28\% increase
in parameters \dots the model is no longer ``lightweight'' compared to the
baseline it improves upon.}

\reply The cost is now reported without softening and measured more carefully than
before (Table~2). Relative to YOLOv11s, the four-level head with MEIS and FRM
carries 21\% more parameters, 2.1 times the operations and 3.5 times the inference
time at 1280\,px. The word ``lightweight'' does not appear as a claim about our
models.

We also found that the batch-size-1 timing convention used in the original
submission understated these costs badly. With one image per batch, two of the
three models were limited by how fast the host could submit work, not by the GPU:
their forward pass took the same time at 1280 and at 1536\,px. Read at face value,
those numbers put the cost of the stride-4 level at 18\% when a GPU-bound
measurement puts it at 89\%. Table~2 therefore reports both batch~1 and batch~8
throughput and marks the host-bound cells.

\point{6. The evaluation is performed only on the car category \dots The +1.8\% AP
gain on large vehicles is not evidence of generalizable small-object detection.}

\reply Agreed. See our response to Reviewer~\#4, point~4. The revised
Section~5.6 reports the car result with sequence-level confidence intervals that
include zero, names the classes UAVDT cannot test, and the abstract states that
small non-vehicle generalization remains untested.

\point{7. The paper reports a gap of 6x its uncertainty, but the uncertainty is
derived from only three random seeds \dots A proper significance test (e.g., a
paired t-test across all validation images) is not performed.}

\reply Section~4.5 is new and describes the testing procedure: a paired bootstrap
over the 548 validation images with 1000 resamples, a paired permutation test on
the same pairing, Bonferroni correction across the ten classes, and sequence-level
resampling for the video data of UAVDT. Every headline contrast is reported with
its confidence interval, and Fig.~6 collects them in a forest plot.

Section~5.4 addresses the seed question directly and separates the two sources of
uncertainty. For the module comparison the paired difference moved from $+0.58$ to
$-0.34$ points between seeds, a standard deviation of 0.65 points, against an
image-level bootstrap standard error of 0.35 --- that is, the seed-to-seed spread
was the larger of the two, and a single seed would have been misleading. The
schedule effect is roughly ten times the seed spread and survives both tests. We
recommend at least three seeds for effects below one point (Section~6.3), and we
report that our own two seeds give only a rough estimate of that variance
(Section~6.4).

\point{8. The paper states that ``both models are evaluated at 1280$\times$1280
resolution,'' but the baseline was only trained at 640px. This gives MEISCF an
unfair advantage \dots}

\reply Correct, and this comparison no longer exists. Every model in the revised
manuscript is trained at its stated input size and evaluated at 1280 and
1536\,px, and each schedule's final training resolution is 1280\,px, so no model
is asked to generalize to an unseen resolution in the headline comparisons
(Sections~4.3--4.4). The only deliberate train--test size mismatch is the 1024\,px
evaluation in Section~5.5, which is used to test a specific hypothesis about the
schedule effect and is labelled as such.

\point{9. Table~2 shows ``YOLOv11 Baseline + Progressive Training'' achieving
50.3\% mAP@50, but it is unclear if this model is trained for the same number of
epochs (350) and with the same hyperparameters. A fair comparison would require a
single-phase training at 1280px for the baseline.}

\reply That experiment is now in the paper and is one of its central results.
YOLOv11s was trained single-phase at 1280\,px under the same budget of up to 350
epochs as its progressive counterpart, and it gained 3.28 points mAP@50 at
1280\,px and 3.97 at 1536\,px over the progressive schedule (Table~5; 95\% CI
$[3.3,\,4.6]$ at 1536\,px, $p<0.002$). Table~3 lists both schedules'
hyperparameters side by side, including the one difference we could not avoid: the
three- and four-level detectors needed different physical batch sizes to fit in
memory, which we disclose in Table~3 and list as a limitation in Section~6.4. We
also note there that this single-phase run stopped early at epoch 345 under a
patience of 200, with its best checkpoint at epoch 145.

\point{10. Learnable weighted fusion of multi-scale features has been explored in
prior work, including EfficientDet's BiFPN and CF-YOLO. The paper's claim that
``simultaneous three-level fusion'' is novel is weak.}

\reply The claim has been withdrawn. See our response to Reviewer~\#4, point~2.

\point{11. Dilated convolutions are typically used to increase receptive fields
\dots but edge detection typically benefits from fine-grained, small-kernel
convolutions. The paper does not explain why dilation rates of 1, 2, and 4 are
optimal for extracting edge information.}

\reply We could not explain it, and the experiment shows there was nothing to
explain: replacing $(1,2,4)$ with $(1,1,1)$ at a fixed parameter count changes
nothing measurable (Section~5.3.1, Table~7). The revised manuscript makes no
optimality claim for the dilation rates and reports this null result instead. The
reviewer's intuition about fine-grained kernels is consistent with what we
measured, and with the sampling argument of Section~3: a module operating after
downsampling cannot recover detail that the stride has already discarded.

\point{12. The paper states that ``ablation confirms balanced weighting
outperforms stronger emphasis,'' but it does not provide the ablation results.}

\reply The sentence has been removed. The residual blend factor $\alpha$ sweep it
referred to belongs to the earlier single-seed configuration and is now reported,
with its actual numbers, in Appendix~A (Table~A.1), together with a statement of
how that configuration differed from the main experiments.

\point{13. The paper inconsistently reports performance at different resolutions
\dots A fair comparison would report FPS at the same resolution as the evaluation.}

\reply Table~2 now reports throughput at 1280 and 1536\,px, which are exactly the
two resolutions at which every model is evaluated. As described in our response to
point~5, we also changed how the measurement is made, because the batch-size-1
convention was measuring the host rather than the model.

\point{14. Reference [26] is cited as ``IEEE Access 10 (2022) 85744-85754,'' but
this appears to be a placeholder or a citation of a different paper \dots}

\reply The reviewer was right that the entry could not be verified. We were unable
to locate it either, and it has been replaced by the paper that actually defines
the MEIS block, Liu et al., \textit{Intelligent and Converged Networks} 6(2),
2025, which is now the reference wherever MEIS-YOLO is discussed. The reference
list was rebuilt for the revised manuscript and contains 30 entries, each cited at
a specific claim in the text. We also checked the author lists and page ranges of
the entries added in revision against the publishers' records, and completed two
that had been abbreviated.

\point{15. The paper states that ``68\% of objects occupy fewer than
32$\times$32 pixels'' and that ``small object features experience 4.2x weaker
activation magnitudes,'' but it does not show the analysis that led to these
numbers.}

\reply Both statements have been dealt with. The activation-magnitude claim was
not supported by any analysis we could reconstruct, and it has been deleted.
The size distribution is now derived rather than asserted. Section~3 gives a
closed-form expression for the smallest object that can receive positive training
samples at a given stride and input size, evaluates it in Table~1, and then
measures the share of VisDrone objects falling below that threshold: at 640\,px,
31.0\% of validation instances are smaller than one stride-8 cell and 5.8\%
smaller than one stride-4 cell; at 1280\,px the same objects give 5.8\% and
0.2\%. These numbers, not the earlier ones, are what the paper's predictions rest
on, and Section~6.1 checks the predictions against the measured results.

% ===========================================================================
\section*{Closing}

The reviewers asked us to isolate our components and to show that the gains
exceeded the noise. We did both, and the answers went against the paper. We have
reported them as we found them, removed the claims they do not support, and
rebuilt the manuscript around the comparison that the evidence does support. We
hope the result is useful, and we thank the reviewers for pushing the work to a
place we would not have reached on our own.

\bigskip

Sincerely, on behalf of all authors,

\bigskip\bigskip

Professor Xinzhong Zhu, corresponding author \\
Moatassim Barakat \\
College of Computer Science and Technology \\
Zhejiang Normal University, Jinhua 321004, China

\end{document}
